% GENERATED FILE. Edit canonical lessons, then run npm run generate:books.
% canonical-source-hash: fnv1a64:1aa95f7433e2e4c4
% canonical-lessons: HI-C400-araksan, HI-C400-barth, HI-C400-pratiksa, HI-C400-jacna, HI-C400-ravana

\chapter{Reservations and Berths}
\label{ch:hi-reservations-and-berths}
\hlchaptermodality{400}

\begin{chapteropening}
\textbf{By the end of this chapter:} \emph{I can say I have a reservation and a berth, wait for a train and have my ticket checked.}

Complete the last lesson of chapter 400: I can say I have a reservation and a berth, wait for a train and have my ticket checked.
\end{chapteropening}

\section[\texorpdfstring{\hi{आरक्षण} (ārakṣaṇ)}{आरक्षण (ārakṣaṇ)}]{\hi{आरक्षण} (ārakṣaṇ) — a reservation (of a seat)}

\label{lesson:HI-C400-araksan}



\begin{quote}
Before the new one: say the Hindi for a whistle, then the Hindi for to slide.
\end{quote}

\subsection*{You'll want to know: \hi{आरक्षण}}
\textbf{\hi{आरक्षण}} — \emph{ārakṣaṇ} — \textquotedblleft{}a reservation (of a seat)\textquotedblright{}.

\begin{quote}
\hi{मेरे} \hi{पास} \hi{इस} \hi{ट्रेन} \hi{में} \hi{आरक्षण} \hi{है।}
\end{quote}

\begin{quote}
\emph{mere pās is ṭren meṁ ārakṣaṇ hai.}
\end{quote}

\begin{quote}
I have a reservation on this train.
\end{quote}

Masculine.

\begin{cousinweb}
\textbf{\hi{आरक्षण}} (\emph{ārakṣaṇ}) is a learned word built from Sanskrit \emph{ā-} and \emph{rakṣaṇa}, \textquotedblleft{}guarding, protecting\textquotedblright{} (Monier-Williams). A reserved seat is one kept safe for you; the railway sense is modern.
\end{cousinweb}

\begin{grammarlens}[title={What you've built}]
One more word for getting around.
\end{grammarlens}

\subsection*{Your turn}
\begin{itemize}
\raggedright
  \item \emph{Say it:} \emph{ārakṣaṇ}
  \item \emph{Say it:} \emph{ārakṣaṇ}, once more
  \item \emph{Say it:} \emph{mere pās is ṭren meṁ ārakṣaṇ hai.} — I have a reservation on this train.
  \item \emph{Recall:} say the Hindi for a whistle, then the Hindi for to slide, then say \emph{ārakṣaṇ} again
\end{itemize}

\begin{tcolorbox}[breakable,colback=teal!4,colframe=teal!35!black,title={Before you move on}]
What does \hi{आरक्षण} mean? (\textquotedblleft{}A reservation (of a seat)\textquotedblright{}.) Say it once more.
\end{tcolorbox}
\section[\texorpdfstring{\hi{बर्थ} (barth)}{बर्थ (barth)}]{\hi{बर्थ} (barth) — a berth (on a train)}

\label{lesson:HI-C400-barth}



\begin{quote}
Before the new one: say the Hindi for to slide, then the Hindi for a reservation (of a seat).
\end{quote}

\subsection*{You'll want to know: \hi{बर्थ}}
\textbf{\hi{बर्थ}} — \emph{barth} — \textquotedblleft{}a berth (on a train)\textquotedblright{}.

\begin{quote}
\hi{मेरी} \hi{बर्थ} \hi{ऊपर} \hi{है} \hi{और} \hi{मीरा} \hi{की} \hi{नीचे।}
\end{quote}

\begin{quote}
\emph{merī barth ūpar hai aur Mīrā kī nīce.}
\end{quote}

\begin{quote}
My berth is the upper one and Mira's is the lower one.
\end{quote}

Feminine.

\begin{cousinweb}
\textbf{\hi{बर्थ}} (\emph{barth}) is borrowed from English \emph{berth}, a sleeping place on a ship or a train; the English th is said with \hi{थ}.
\end{cousinweb}

\begin{grammarlens}[title={What you've built}]
One more word for getting around.
\end{grammarlens}

\subsection*{Your turn}
\begin{itemize}
\raggedright
  \item \emph{Say it:} \emph{barth}
  \item \emph{Say it:} \emph{barth}, once more
  \item \emph{Say it:} \emph{merī barth ūpar hai aur Mīrā kī nīce.} — My berth is the upper one and Mira's is the lower one.
  \item \emph{Recall:} say the Hindi for to slide, then the Hindi for a reservation (of a seat), then say \emph{barth} again
\end{itemize}

\begin{tcolorbox}[breakable,colback=teal!4,colframe=teal!35!black,title={Before you move on}]
What does \hi{बर्थ} mean? (\textquotedblleft{}A berth (on a train)\textquotedblright{}.) Say it once more.
\end{tcolorbox}
\section[\texorpdfstring{\hi{प्रतीक्षा} (pratīkṣā)}{प्रतीक्षा (pratīkṣā)}]{\hi{प्रतीक्षा} (pratīkṣā) — waiting; a wait (formal)}

\label{lesson:HI-C400-pratiksa}



\begin{quote}
Before the new one: say the Hindi for a reservation (of a seat), then the Hindi for a berth.
\end{quote}

\subsection*{You'll want to know: \hi{प्रतीक्षा}}
\textbf{\hi{प्रतीक्षा}} — \emph{pratīkṣā} — \textquotedblleft{}waiting; a wait (formal)\textquotedblright{}.

\begin{quote}
\hi{हम} \hi{दो} \hi{घंटे} \hi{से} \hi{ट्रेन} \hi{की} \hi{प्रतीक्षा} \hi{कर} \hi{रहे} \hi{हैं।}
\end{quote}

\begin{quote}
\emph{ham do ghaṇṭe se ṭren kī pratīkṣā kar rahe haiṁ.}
\end{quote}

\begin{quote}
We have been waiting for the train for two hours.
\end{quote}

Feminine.

Station signs read \hi{प्रतीक्षालय}, the waiting room; a ticket on the waiting list is on the \hi{प्रतीक्षा} \hi{सूची}.

\begin{cousinweb}
\textbf{\hi{प्रतीक्षा}} (\emph{pratīkṣā}) is a tatsama, the Sanskrit word \emph{pratīkṣā}, \textquotedblleft{}expectation, looking out for\textquotedblright{} (Monier-Williams), kept unchanged. It belongs to signs and announcements; in talk people say \hi{इंतज़ार}.
\end{cousinweb}

\begin{grammarlens}[title={What you've built}]
One more word for getting around.
\end{grammarlens}

\subsection*{Your turn}
\begin{itemize}
\raggedright
  \item \emph{Say it:} \emph{pratīkṣā}
  \item \emph{Say it:} \emph{pratīkṣā}, once more
  \item \emph{Say it:} \emph{ham do ghaṇṭe se ṭren kī pratīkṣā kar rahe haiṁ.} — We have been waiting for the train for two hours.
  \item \emph{Recall:} say the Hindi for a reservation (of a seat), then the Hindi for a berth, then say \emph{pratīkṣā} again
\end{itemize}

\begin{tcolorbox}[breakable,colback=teal!4,colframe=teal!35!black,title={Before you move on}]
What does \hi{प्रतीक्षा} mean? (\textquotedblleft{}Waiting; a wait (formal)\textquotedblright{}.) Say it once more.
\end{tcolorbox}
\section[\texorpdfstring{\hi{जाँचना} (jā̃cnā)}{जाँचना (jā̃cnā)}]{\hi{जाँचना} (jā̃cnā) — to check; to examine}

\label{lesson:HI-C400-jacna}



\begin{quote}
Before the new one: say the Hindi for a berth, then the Hindi for waiting.
\end{quote}

\subsection*{You'll want to know: \hi{जाँचना}}
\textbf{\hi{जाँचना}} — \emph{jā̃cnā} — \textquotedblleft{}to check; to examine\textquotedblright{}.

\begin{quote}
\hi{एक} \hi{आदमी} \hi{ने} \hi{हमारे} \hi{टिकट} \hi{जाँचे।}
\end{quote}

\begin{quote}
\emph{ek ādmī ne hamāre ṭikaṭ jā̃ce.}
\end{quote}

\begin{quote}
A man checked our tickets.
\end{quote}

Past, with \hi{ने} (he or she …): \hi{उसने} \hi{टिकट} \hi{जाँचा}, \emph{usne ṭikaṭ jā̃cā}.

\begin{cousinweb}
\textbf{\hi{जाँचना}} (\emph{jā̃cnā}) is a tadbhava of Sanskrit \emph{yācyate}, \textquotedblleft{}is asked\textquotedblright{} (CDIAL 10449), which lists Hindi \emph{jā̃cnā}, \textquotedblleft{}to test\textquotedblright{}. Asking turned into questioning a thing until you know it is right.
\end{cousinweb}

\begin{grammarlens}[title={What you've built}]
One more word for getting around.
\end{grammarlens}

\subsection*{Your turn}
\begin{itemize}
\raggedright
  \item \emph{Say it:} \emph{jā̃cnā}
  \item \emph{Say it:} \emph{jā̃cnā}, once more
  \item \emph{Say it:} \emph{ek ādmī ne hamāre ṭikaṭ jā̃ce.} — A man checked our tickets.
  \item \emph{Recall:} say the Hindi for a berth, then the Hindi for waiting, then say \emph{jā̃cnā} again
\end{itemize}

\begin{tcolorbox}[breakable,colback=teal!4,colframe=teal!35!black,title={Before you move on}]
What does \hi{जाँचना} mean? (\textquotedblleft{}To check; to examine\textquotedblright{}.) Say it once more.
\end{tcolorbox}
\section[\texorpdfstring{\hi{रवाना} (ravānā)}{रवाना (ravānā)}]{\hi{रवाना} (ravānā) — setting off; on its way}

\label{lesson:HI-C400-ravana}



\begin{quote}
Before the new one: say the Hindi for waiting, then the Hindi for to check.
\end{quote}

\subsection*{You'll want to know: \hi{रवाना}}
\textbf{\hi{रवाना}} — \emph{ravānā} — \textquotedblleft{}setting off; on its way\textquotedblright{}.

\begin{quote}
\hi{ट्रेन} \hi{सुबह} \hi{छह} \hi{बजे} \hi{रवाना} \hi{हुई।}
\end{quote}

\begin{quote}
\emph{ṭren subah chah baje ravānā huī.}
\end{quote}

\begin{quote}
The train set off at six in the morning.
\end{quote}

\hi{रवाना} \hi{होना}, to set off; \hi{रवाना} \hi{करना}, to send off.

\begin{cousinweb}
\textbf{\hi{रवाना}} (\emph{ravānā}) is usually traced to Persian \emph{ravāna}, \textquotedblleft{}going, flowing, setting out\textquotedblright{}. It does not change for gender: \hi{ट्रेन} \hi{रवाना} \hi{हुई}, the train left.
\end{cousinweb}

\begin{grammarlens}[title={What you've built}]
One more word for getting around.
\end{grammarlens}

\subsection*{Your turn}
\begin{itemize}
\raggedright
  \item \emph{Say it:} \emph{ravānā}
  \item \emph{Say it:} \emph{ravānā}, once more
  \item \emph{Say it:} \emph{ṭren subah chah baje ravānā huī.} — The train set off at six in the morning.
  \item \emph{Recall:} say the Hindi for waiting, then the Hindi for to check, then say \emph{ravānā} again
\end{itemize}

\begin{tcolorbox}[breakable,colback=teal!4,colframe=teal!35!black,title={Before you move on}]
What does \hi{रवाना} mean? (\textquotedblleft{}Setting off; on its way\textquotedblright{}.) Say it once more.
\end{tcolorbox}
